\documentclass[14pt,a4paper,oneside]{extbook}
\usepackage[utf8]{vietnam}
\usepackage{amsmath,amssymb,mathrsfs,fancyhdr,enumerate,multirow,makecell}
\usepackage{tkz-euclide,tikz-3dplot,tikz,tkz-tab}
\usetikzlibrary{arrows.meta,calc,intersections,angles,quotes,backgrounds,shapes.geometric,patterns,through,shadows,positioning}
\usepackage{yhmath}
\usepackage[top=0.7cm, bottom=0.7cm, left=1.0cm, right=1.0cm, headheight=16pt, headsep=3mm, footskip=5mm]{geometry}
\usepackage[hidelinks,unicode]{hyperref}
\usepackage{qrcode}
\usepackage[many]{tcolorbox}
\usepackage{fontawesome}

\pagestyle{fancy}
\fancyhf{}
\fancyhead[L]{\footnotesize\bfseries\color{teal!85!black}\faBook~SỔ TAY KIẾN THỨC TOÁN 4}
\fancyhead[R]{\footnotesize\bfseries\color{blue!85!black}Lớp Toán Thầy Lĩnh -- loptoanthaylinh.vn}
\fancyfoot[C]{\footnotesize\bfseries\color{gray!80!black}-- Sổ tay kiến thức: Bài 70. Ôn tập các phép tính với phân số --}
\renewcommand{\headrulewidth}{0.4pt}

\begin{document}

% --- BANNER TIÊU ĐỀ BÀI HỌC ---
\noindent
\begin{tcolorbox}[
  enhanced,
  colback=teal!8!white,
  colframe=teal!80!black,
  arc=3mm, boxrule=1pt,
  left=3mm, right=3mm, top=2mm, bottom=2mm,
  halign=center
]
  {\Large\bfseries\color{teal!90!black}\faBookmark\ BÀI 70: ÔN TẬP CÁC PHÉP TÍNH VỚI PHÂN SỐ}\\[1.5pt]
  {\footnotesize\itshape\color{black!75} Bảng đối chiếu 4 phép tính, quy tắc vàng tính toán và tìm phân số của một số trên 1 trang A4}
\end{tcolorbox}

\vspace{-2.5mm}

\noindent
\begin{minipage}[t]{0.49\textwidth}\vspace{0pt}
\begin{tcolorbox}[
  enhanced,
  colback=blue!4!white,
  colframe=blue!75!black,
  arc=2mm, boxrule=0.8pt,
  left=2mm, right=2mm, top=1.8mm, bottom=1.8mm,
  title={\footnotesize\bfseries\color{white}\faPlusSquare\ 1. Phép Cộng và Phép Trừ}
]
{\footnotesize
\begin{itemize}\setlength{\itemsep}{1pt}\setlength{\leftskip}{-2.5mm}
  \item \textbf{Cùng mẫu số:} Cộng/trừ tử số, \textbf{giữ nguyên mẫu số}:
  $$\dfrac{a}{m} \pm \dfrac{b}{m} = \dfrac{a \pm b}{m}$$
  \item \textbf{Khác mẫu số:} Bắt buộc phải \textbf{Quy đồng mẫu số} trước, sau đó cộng/trừ như cùng mẫu.
  \item Tuyệt đối \textbf{không} cộng/trừ mẫu với mẫu!
\end{itemize}
}
\end{tcolorbox}
\end{minipage}
\hfill
\begin{minipage}[t]{0.49\textwidth}\vspace{0pt}
\begin{tcolorbox}[
  enhanced,
  colback=teal!4!white,
  colframe=teal!75!black,
  arc=2mm, boxrule=0.8pt,
  left=2mm, right=2mm, top=1.8mm, bottom=1.8mm,
  title={\footnotesize\bfseries\color{white}\faTimesCircle\ 2. Phép Nhân và Phép Chia}
]
{\footnotesize
\begin{itemize}\setlength{\itemsep}{1pt}\setlength{\leftskip}{-2.5mm}
  \item \textbf{Phép nhân:} Lấy tử nhân tử, mẫu nhân mẫu:
  $$\dfrac{a}{b} \times \dfrac{c}{d} = \dfrac{a \times c}{b \times d}$$
  (Nên triệt tiêu rút gọn chéo trước khi nhân).
  \item \textbf{Phép chia:} Nhân với phân số đảo ngược:
  $$\dfrac{a}{b} : \dfrac{c}{d} = \dfrac{a}{b} \times \dfrac{d}{c} = \dfrac{a \times d}{b \times c}$$
\end{itemize}
}
\end{tcolorbox}
\end{minipage}

\vspace{-2.5mm}
% --- KHỐI 3: SƠ ĐỒ TƯ DUY TIKZ ---
\noindent
\begin{tcolorbox}[
  enhanced,
  colback=orange!4!white,
  colframe=orange!80!black,
  arc=2mm, boxrule=0.8pt,
  left=2mm, right=2mm, top=1.5mm, bottom=1.5mm,
  title={\footnotesize\bfseries\color{white}\faSitemap\ SƠ ĐỒ TƯ DUY: MA TRẬN 4 PHÉP TÍNH PHÂN SỐ}
]
\centering
\begin{tikzpicture}[
  box/.style={draw, rounded corners=2pt, align=center, font=\scriptsize\bfseries, inner sep=3pt, minimum width=3.4cm, minimum height=0.95cm}
]
  \node[box, fill=blue!12, draw=blue!80!black] at (1.8, 0.6) {PHÉP CỘNG (+)\\Cùng mẫu: Cộng tử, giữ mẫu\\Khác mẫu: Quy đồng};

  \node[box, fill=teal!12, draw=teal!80!black] at (6.0, 0.6) {PHÉP TRỪ (--)\\Cùng mẫu: Trừ tử, giữ mẫu\\Khác mẫu: Quy đồng};

  \node[box, fill=orange!12, draw=orange!80!black] at (1.8, -0.6) {PHÉP NHÂN ($\times$)\\Tử $\times$ Tử, Mẫu $\times$ Mẫu\\Triệt tiêu trước khi nhân};

  \node[box, fill=green!15, draw=green!70!black] at (6.0, -0.6) {PHÉP CHIA (:)\\Nhân với phân số đảo ngược\\$\frac{a}{b} \times \frac{d}{c}$};

  \node[draw=purple!80!black, fill=purple!15, rounded corners=3pt, font=\scriptsize\bfseries, align=center, text width=2.5cm] at (10.1, 0) {KẾT QUẢ\\LUÔN ĐƯA VỀ\\TỐI GIẢN!};
  \draw[-{Stealth}, purple!80!black, thick] (7.8, 0) -- (8.7, 0);
\end{tikzpicture}
\end{tcolorbox}

\vspace{-2.5mm}
% --- KHỐI 4: VÍ DỤ & MẸO TRÁNH BẪY ---
\noindent
\begin{minipage}[t]{0.49\textwidth}\vspace{0pt}
\begin{tcolorbox}[
  enhanced,
  colback=green!4!white,
  colframe=green!60!black,
  arc=2mm, boxrule=0.8pt,
  left=2mm, right=2mm, top=1.8mm, bottom=1.8mm,
  title={\footnotesize\bfseries\color{white}\faCheckCircle\ Ví dụ thực hành chuẩn mực}
]
{\footnotesize
\textbf{Tính giá trị biểu thức:}\\
$C = \left(\dfrac{1}{2} + \dfrac{1}{3}\right) \times \dfrac{6}{5}$\\
Trong ngoặc: $\dfrac{1}{2} + \dfrac{1}{3} = \dfrac{3}{6} + \dfrac{2}{6} = \dfrac{5}{6}$.\\
Nhân: $\dfrac{5}{6} \times \dfrac{6}{5} = \dfrac{5 \times 6}{6 \times 5} = \mathbf{1}$.
}
\end{tcolorbox}
\end{minipage}
\hfill
\begin{minipage}[t]{0.49\textwidth}\vspace{0pt}
\begin{tcolorbox}[
  enhanced,
  colback=red!4!white,
  colframe=red!70!black,
  arc=2mm, boxrule=0.8pt,
  left=2mm, right=2mm, top=1.8mm, bottom=1.8mm,
  title={\footnotesize\bfseries\color{white}\faExclamationTriangle\ Mẹo tránh bẫy trắc nghiệm}
]
{\footnotesize
\begin{itemize}\setlength{\itemsep}{1.5pt}\setlength{\leftskip}{-2.5mm}
  \item \textbf{Bẫy nhầm giữa Cộng và Nhân:}\\
  Cộng phân số KHÔNG lấy mẫu cộng mẫu. Nhân phân số PHẢI lấy mẫu nhân mẫu!
  \item \textbf{Bẫy chia 1 cho phân số:}\\
  $1 : \frac{a}{b} = \frac{b}{a}$ (bằng chính phân số đảo ngược).
\end{itemize}
}
\end{tcolorbox}
\end{minipage}

\end{document}
